% coalition.tex -- Section 7: Coalition Formation

\section{Coalition Formation}
\label{sec:coalition}

Coalition formation extends N-player games with a structured negotiation phase,
enabling agents to form binding or non-binding agreements before acting. This
tests whether agents can build trust, coordinate collectively, and handle
betrayal---capabilities directly relevant to multi-agent alignment.

Coalition formation tests whether agents can build trust, coordinate
collectively, and handle betrayal---capabilities central to multi-agent
alignment in deployment settings where agents must negotiate, form
agreements, and decide whether to honor commitments.

\subsection{Two-Phase Protocol}
\label{sec:coalition-protocol}

Each round of a coalition game proceeds in two phases:

\begin{enumerate}[nosep]
    \item \textbf{Negotiate Phase}: players exchange \texttt{CoalitionProposal}
          objects specifying target members, an agreed joint action, and an optional
          side payment. Other players accept or reject via their coalition strategy.
    \item \textbf{Act Phase}: players choose actions independently. Coalition
          members may honor or violate their agreed action, subject to the
          enforcement mode. The environment records defectors by comparing
          each member's chosen action against \texttt{coalition.agreed\_action}.
\end{enumerate}

\subsection{Enforcement Modes}
\label{sec:enforcement}

Three enforcement modes model different institutional settings:

\begin{description}[nosep]
    \item[Cheap Talk.] Agreements are non-binding. Players may freely
        deviate with no payoff modification.
    \item[Penalty.] Defectors pay a proportional penalty:
        $p_d = u_d \cdot \frac{\alpha}{\beta}$ (default $\frac{1}{2}$),
        yielding adjusted payoff $\hat{u}_d = u_d - p_d = \frac{u_d}{2}$.
    \item[Binding.] Actions are overridden at the environment level:
        coalition members' actions are locked to \texttt{agreed\_action}
        before payoff computation, preventing defection entirely.
\end{description}

\subsection{Side Payments and Defection Detection}
\label{sec:side-payments}

The coalition framework supports side payments and defection detection.

\paragraph{Side payments.} The proposer (first coalition member) pays each
other member a fixed amount $s$ (default $0$):
\begin{align}
\hat{u}_{\text{proposer}} &= u_{\text{proposer}} - s \cdot (|C| - 1) \\
\hat{u}_{m} &= u_m + s \quad \forall m \in C \setminus \{\text{proposer}\}
\end{align}

\paragraph{Defection detection.} After the Act phase, the environment
compares each coalition member's action against \texttt{agreed\_action}.
Any mismatch adds the member to the defectors list, which is visible to
all players in subsequent rounds, enabling retaliatory strategies.

\subsection{Coalition Strategies}
\label{sec:coalition-strategies}

Built-in coalition strategies model different agent dispositions:

\begin{description}[nosep]
    \item[Loyal.] Accepts all proposals; always plays \texttt{agreed\_action}.
        Models unconditional cooperators.
    \item[Betrayer.] Accepts all proposals but plays the first \emph{non-agreed}
        action, maximizing exploitation of trusting partners.
    \item[Conditional.] Honors agreements unless any other member defected
        in the previous round (checked via \texttt{coalition\_history[-1].defectors}).
        Models reciprocal enforcement.
    \item[Random.] Accepts/rejects proposals with probability $\frac{1}{2}$;
        chooses actions uniformly at random.
\end{description}

\subsection{Coalition Game Library}
\label{sec:coalition-games}

Seven games are specifically designed for coalition dynamics:

\begin{itemize}[nosep]
    \item \textbf{Cartel Formation} --- firms collude on pricing vs.\ compete
    \item \textbf{Alliance Game} --- mutual defense pacts with defection risk
    \item \textbf{Voting Blocs} --- coordinated voting for collective benefit
    \item \textbf{Ostracism Game} --- coalition excludes non-cooperators
    \item \textbf{Resource Trading} --- bilateral exchange within coalitions
    \item \textbf{Rule Voting} --- coalitions propose and vote on rule changes
    \item \textbf{Commons Management} --- coalition-based resource governance
\end{itemize}

These games combine the base game mechanics with coalition negotiation,
creating rich strategic environments where agents must balance individual
payoff maximization against coalition stability and fairness.
